% s-123-p123-m002; original PDF 123
\Aphorism{s-123-p123-m002}{158}
Um den Weg nicht bloss als Weg zu gehen, lügt man ein Ziel hinzu; der Zweck ist Reflex der Mittel. Mittel und Wege nämlich haben kein gutes Gewissen, sondern erwerben es erst auf dem Umweg über Zwecke und Ziele; dem Menschen ohne Teleologie wird es zu Muthe wie einem epikureischen Gotte in seiner Zwischenwelt: selig, aber zweifelhaft, ob er zum Leben berechtigt sei. Dass das Verschwenden von Mitteln, auch ohne Zweck, gesund ist, dass man das Recht hat zu brennen, auch ohne die Öfen der Utilität zu heizen, ist uns noch neu: als sesshafte Ackerbauer und \Foreign{glebae adscripti} glauben wir nicht daran, dass Unterwegs auch ein Ziel sei. Der Nomadenzustand, die grosse Freiheit und Freizügigkeit, soll erst noch kommen, wozu allererst erforderlich ist, die Weiblein dafür zu gewinnen. Denn das Weib ist von der Nothwendigkeit, Ziele zu haben, d. h. etwas Nahes, Bestimmtes, Greifbares zu wollen, axiomatisch überzeugt: und sein eigenes Ziel ist der Mann, der solchermassen ziel- und zweckgemäss existirt. Ein tief philisterhafter Grundinstinct, mit dem es den Mann, der unter Umständen Nichtphilister sein will, immer wieder zurückbildet; denn wenn er das Weib gewinnen will, ist er auch verurtheilt die Waffen zu führen, mit denen der Kampf um das Weib ausgefochten wird. Das Weib lässt sich Zukünfte zur Auswahl vorlegen, und es ist wahrhaftig nicht die fernere, zukünftigere Zukunft, in die es Lust hätte hineinzuspringen.
\par

